\section{Producto operativo: catálogo web de construcciones tradicionales}

Como ocurre en varias ICR orientadas por la conservación del patrimonio construido, esta investigación no concluye únicamente con una interpretación teórica. El producto operativo será un catálogo web de consulta que permita reunir fichas, fotografías, descripciones, mapas generalizados y criterios de documentación sobre construcciones tradicionales hñähñu.

La decisión de plantear el producto como página web responde a una necesidad práctica: el registro no debe quedar cerrado en cuatro casos ni limitado al documento académico. La plataforma permitirá iniciar con las construcciones estudiadas en esta ICR y, posteriormente, incorporar nuevos sistemas, comunidades o registros de campo.

\subsection{Objetivo del catálogo web}

El catálogo tendrá tres funciones principales:

\begin{enumerate}
	\item \textbf{Consultar}: presentar información clara sobre cada construcción documentada, con fotografías, descripción del sistema y ubicación generalizada.
	\item \textbf{Comparar}: permitir la lectura entre sistemas constructivos, materiales, usos y estados de conservación.
	\item \textbf{Ampliar}: ofrecer una estructura de datos que pueda crecer con nuevos registros sin rehacer el sistema de clasificación.
\end{enumerate}

De esta manera, la web funcionará como una herramienta de documentación y divulgación patrimonial, no como un repositorio cerrado ni como una propuesta de intervención constructiva.

\subsection{Usuarios previstos}

El producto se dirige a cuatro tipos de usuarios:

\begin{itemize}
	\item habitantes o integrantes de las comunidades que deseen consultar el registro;
	\item estudiantes e investigadores interesados en arquitectura vernácula y patrimonio biocultural;
	\item docentes o asesores que requieran revisar la evidencia organizada;
	\item futuros colaboradores que puedan ampliar el catálogo con nuevas construcciones.
\end{itemize}

La interfaz deberá privilegiar una lectura sencilla: fichas visibles, fotografías claras, filtros básicos y mapas que no expongan información sensible.

\subsection{Estructura de información}

Cada registro del catálogo corresponderá a una construcción. La ficha web deberá organizarse con los siguientes campos:

\begin{table}[ht]
	\centering
	\caption{Campos mínimos para la ficha web}
	\vspace{0.25cm}
	\label{tab:campos_ficha_web}
	\begin{tabular}{@{}L{0.28\linewidth}L{0.62\linewidth}@{}}
		\toprule
		\textbf{Campo} & \textbf{Contenido} \\
		\midrule
		Clave & Identificador anónimo de la construcción. \\
		Sistema constructivo & Penca de maguey, piedra local, bajareque, órganos u otro sistema añadido. \\
		Comunidad & Nombre público o versión generalizada, según criterios de privacidad. \\
		Uso & Cocina, bodega, corral, resguardo, cerramiento u otro uso registrado. \\
		Materiales & Tierra, piedra, penca, lechuguilla, madera, órgano, lámina u otros. \\
		Estado & Descripción breve de permanencias, deterioros y sustituciones. \\
		Relato asociado & Síntesis del testimonio, sin datos personales identificables. \\
		Fotografías & Imágenes seleccionadas y autorizadas para consulta. \\
		Mapa & Localización generalizada o referencia territorial no exacta. \\
		Observaciones & Notas sobre cambios, mantenimiento o dudas pendientes. \\
		\bottomrule
	\end{tabular}
\end{table}

Esta estructura permite que la ficha sea legible en la tesis, en una página web y en una posible versión impresa. También facilita exportar los datos a formatos tabulares si en el futuro se requiere continuar el inventario.

\subsection{Organización del sitio}

La página web se propone con cuatro secciones principales:

\begin{enumerate}
	\item \textbf{Inicio}: breve explicación del proyecto, alcance territorial y sentido del catálogo.
	\item \textbf{Sistemas constructivos}: descripción general de cada sistema identificado, con fotografías representativas.
	\item \textbf{Catálogo}: fichas individuales de las construcciones documentadas, con filtros por sistema, material, comunidad y estado.
	\item \textbf{Mapa}: visualización generalizada de las construcciones y comunidades, cuidando la privacidad de ubicaciones exactas.
\end{enumerate}

De forma complementaria, podrá agregarse una sección de \emph{metodología} para explicar cómo se registraron los casos y una sección de \emph{créditos y consentimiento} para transparentar el uso de imágenes y testimonios.

\subsection{Crecimiento del catálogo}

El catálogo deberá diseñarse para crecer. Para ello, cada nuevo registro deberá seguir la misma ficha base y asignarse a un sistema constructivo existente o a una nueva categoría justificada. El crecimiento se plantea en tres niveles:

\begin{enumerate}
	\item \textbf{Nuevas construcciones}: incorporación de más casos dentro de El Deca y El Buena.
	\item \textbf{Nuevas comunidades}: ampliación a otras localidades del Valle del Mezquital.
	\item \textbf{Nuevos sistemas}: registro de variantes constructivas que no aparezcan en el corpus inicial.
\end{enumerate}

Esta lógica evita que el producto dependa únicamente de los cuatro casos iniciales y permite que la ICR funcione como primera fase de un archivo patrimonial más amplio.

\subsection{Criterios de documentación y conservación preventiva}

A partir del diagnóstico se formularán criterios básicos para documentar y conservar estas construcciones sin forzar una propuesta de restauración. Los criterios preliminares son:

\begin{enumerate}
	\item \textbf{Registrar antes de intervenir}: documentar fotografía, ubicación general, uso y relato antes de cualquier modificación.
	\item \textbf{Reconocer el sistema completo}: no separar muros, cubierta, piso, patio, vegetación y usos cotidianos como si fueran elementos aislados.
	\item \textbf{Priorizar compatibilidad material}: evitar sustituciones que generen incompatibilidades físicas o alteren de forma irreversible la lectura constructiva.
	\item \textbf{Registrar cambios como parte de la historia}: distinguir entre deterioro, mantenimiento, adaptación y sustitución.
	\item \textbf{Proteger la información sensible}: generalizar ubicaciones y anonimizar construcciones cuando la publicación pueda afectar a propietarios o habitantes.
	\item \textbf{Incluir la voz local}: asociar cada ficha con relatos de uso, memoria o mantenimiento para no reducir la construcción a su materialidad.
\end{enumerate}

Estos criterios serán la base ética y metodológica del catálogo web. Su función es garantizar que la ampliación del registro no convierta las construcciones en objetos aislados, sino que mantenga visible la relación entre materialidad, territorio y saber local.

\subsection{Privacidad y publicación}

La versión pública del catálogo no deberá publicar coordenadas exactas ni datos personales de propietarios o entrevistados. Cuando una fotografía permita identificar a una persona, un predio o una situación sensible, se usará solo con autorización explícita o se reservará para consulta interna.

El catálogo podrá manejar dos niveles de información: una versión pública con ubicación generalizada y una versión interna con datos completos para revisión académica. Esta distinción permite aprovechar la precisión del trabajo cartográfico sin exponer indebidamente a las comunidades.
